% !TeX root = ../analysis1.tex
\chapter{Limits of Functions and Continuity}\label{chap:analysis-continuity}

\section{Limits and one-sided limits}
A function limit describes values near a point. Its value at that point
can be different, or the function may not be defined there at all.

\begin{definition}[Limit at an interior point]
Suppose $f$ is defined on an open interval about $c$, except possibly at $c$.
We write $\lim_{x\to c}f(x)=L$ if
\[
\forall\varepsilon>0\ \exists\delta>0:\quad
0<|x-c|<\delta\ \Longrightarrow\ |f(x)-L|<\varepsilon.
\]
The condition $0<|x-c|$ excludes the point $c$ itself.
\end{definition}

\begin{definition}[One-sided limits]
The right-hand limit $\lim_{x\to a^+}f(x)=L$ means that for every
$\varepsilon>0$ there is $\delta>0$ such that
\[
a<x<a+\delta\ \Longrightarrow\ |f(x)-L|<\varepsilon,
\]
for $x$ in the domain. The left-hand limit at $b$ uses $b-\delta<x<b$.
We also write $x\downarrow a$ and $x\uparrow b$, respectively.
\end{definition}
At an interior point a two-sided limit exists with value $L$ exactly when
both one-sided limits exist and equal $L$.

\begin{example}[Different one-sided limits]
Let
\[
f(x)=\begin{cases}2x,&x<1,\\1,&x=1,\\4-x,&x>1.\end{cases}
\]
Then $\lim_{x\to1^-}f(x)=2$ and $\lim_{x\to1^+}f(x)=3$, so the two-sided
limit does not exist. For the left limit,
$|f(x)-2|=2|x-1|$, so choose $\delta=\varepsilon/2$.
For the right limit, $|f(x)-3|=|x-1|$, so choose $\delta=\varepsilon$.
The assigned value $f(1)=1$ does not affect either limit.
\end{example}

\begin{example}[Oscillation controlled by a small factor]
For $x\ne0$, let $f(x)=1+x^2\cos(1/x)$. Since
\[
|f(x)-1|\leq x^2,
\]
choosing $\delta=\sqrt\varepsilon$ proves $\lim_{x\to0}f(x)=1$.
No value of $f(0)$ was needed.
\end{example}

\section{Continuity}
\begin{definition}[Continuity on a domain]
Let $S\subseteq\mathbb R$ and $c\in S$. A function $f:S\to\mathbb R$
is continuous at $c$ if
\[
\forall\varepsilon>0\ \exists\delta>0\ \forall x\in S:\quad
|x-c|<\delta\ \Longrightarrow\ |f(x)-f(c)|<\varepsilon.
\]
It is continuous on $S$, written $f\in C(S)$, if it is continuous at every
point of $S$. Otherwise it is discontinuous at any point where the condition
fails.
\end{definition}
At an interior point of an interval, continuity is equivalent to
$\lim_{x\to c}f(x)=f(c)$. On $[a,b]$, the endpoint conditions are
\[
\lim_{x\to a^+}f(x)=f(a),\qquad
\lim_{x\to b^-}f(x)=f(b).
\]
Only points in the domain are tested.

\begin{example}[A removable discontinuity]
Extend $1+x^2\cos(1/x)$ to $0$ by setting $f(0)=\gamma$.
The extension is continuous at $0$ if and only if $\gamma=1$.
\end{example}

\begin{example}[Thomae's function]
On $(0,1)$ define
\[
T(x)=\begin{cases}
0,&x\notin\mathbb Q,\\
1/q,&x=p/q\text{ in lowest terms},\ q>0.
\end{cases}
\]
Then $\lim_{x\to c}T(x)=0$ for every $c\in(0,1)$. It is continuous at
every irrational point and discontinuous at every rational point.

\textit{Supplementary explanation.} Given $\varepsilon>0$, choose an
integer $Q>1/\varepsilon$. There are only finitely many reduced fractions
in $(0,1)$ with denominator at most $Q$. A sufficiently small punctured
neighbourhood of $c$ avoids all of them, other than possibly $c$ itself.
At any other rational point in that neighbourhood, $q>Q$, so
$T(x)<\varepsilon$; at an irrational point $T(x)=0$.
Therefore the punctured limit is zero. This equals $T(c)$ exactly when
$c$ is irrational.
\end{example}

\section{Sequential criteria and algebra}
\begin{theorem}[Sequential criterion for limits]\label{thm:analysis-sequential}
Under the domain assumptions for an interior limit,
$\lim_{x\to c}f(x)=L$ if and only if every sequence $(x_n)$ in the domain
with $x_n\ne c$ and $x_n\to c$ satisfies $f(x_n)\to L$.
For a one-sided limit, require the sequence to approach from that side.
\end{theorem}
\begin{proof}
Suppose the limit exists. Given $\varepsilon>0$, choose a corresponding
$\delta>0$. Since $x_n\to c$, eventually $0<|x_n-c|<\delta$, and therefore
$|f(x_n)-L|<\varepsilon$.

Conversely, if the limit condition fails, there is $\varepsilon_0>0$ such
that for every $\delta>0$ there is a domain point $x$ satisfying
$0<|x-c|<\delta$ and $|f(x)-L|\geq\varepsilon_0$.
Choose such a point $x_n$ for $\delta=1/n$. Then $x_n\to c$, but
$f(x_n)\not\to L$, contradicting the sequential hypothesis.
Restrict the selected points to the appropriate side for one-sided limits.
\end{proof}

\begin{corollary}[Sequential continuity]
A function $f:S\to\mathbb R$ is continuous at $c\in S$ if and only if
for every sequence $x_n\in S$ tending to $c$, we have $f(x_n)\to f(c)$.
\end{corollary}
The proof is the same, without requiring $x_n\ne c$.

\begin{theorem}[Algebra of function limits]
If $f(x)\to L$ and $g(x)\to M$ as $x\to c$, then
\[
f(x)+g(x)\to L+M,\qquad f(x)g(x)\to LM.
\]
If $M\ne0$, then $g$ is nonzero sufficiently near $c$ and
$f(x)/g(x)\to L/M$. The analogous rules hold for one-sided limits.
In particular, sums and products of functions continuous at $c$ are
continuous there; their quotient is continuous if $g(c)\ne0$.
\end{theorem}
\begin{proof}
Apply the sequence limit rules to $f(x_n)$ and $g(x_n)$ for every sequence
approaching $c$, then use Theorem~\ref{thm:analysis-sequential}.
Nonvanishing follows by making $|g(x)-M|<|M|/2$.
\end{proof}

\begin{theorem}[Sandwich theorem for functions]
If $f(x)\leq g(x)\leq h(x)$ sufficiently near $c$, and both bounding
functions tend to $L$ as $x\to c$, then $g(x)\to L$.
Continuity of $g$ is not assumed.
\end{theorem}
\begin{proof}
Apply the sequence sandwich theorem to $f(x_n),g(x_n),h(x_n)$ for any
sequence approaching $c$, or use the same $L\pm\varepsilon$ bounds directly.
\end{proof}
\begin{example}
For $x>0$, $-\sqrt x\leq\sqrt x\cos(1/x)\leq\sqrt x$. Both bounds tend
to zero as $x\to0^+$, so $\sqrt x\cos(1/x)\to0$.
\end{example}

\begin{example}[Standard continuous functions]
Constant functions and $f(x)=x$ are continuous directly from the definition.
The algebra rules show that every polynomial is continuous, and every
rational function $p(x)/q(x)$ is continuous wherever $q(x)\ne0$.
For instance $(x^2-2x+7)/(x^2+1)$ is continuous on $\mathbb R$.
\end{example}

\begin{theorem}[Composition]
If $g$ is continuous at $c$ and $f$ is continuous at $g(c)$, with the
composition defined near $c$ in its domain, then $f\circ g$ is continuous
at $c$.
\end{theorem}
\begin{proof}
For $x_n\to c$, continuity gives $g(x_n)\to g(c)$, then
$f(g(x_n))\to f(g(c))$.
\end{proof}

\section{Continuity of the exponential}
Recall the definition $\exp(x)=\sum_{n=0}^{\infty}x^n/n!$ and the addition
law from Chapter~\ref{chap:analysis-series}. We already know positivity
and $\exp(-x)=1/\exp(x)$.

\begin{lemma}[Bounds near zero]\label{lem:analysis-exp-bounds}
For $0\leq x<1$,
\[
1+x\leq\exp(x)\leq\sum_{n=0}^{\infty}x^n=\frac1{1-x}.
\]
For $-1<x<0$, writing $t=-x\in(0,1)$ and taking reciprocals gives
\[
1+x\leq\exp(x)\leq\frac1{1-x}.
\]
Hence the same bounds hold throughout $|x|<1$.
\end{lemma}
\begin{proof}
For nonnegative $x$, use positivity of the exponential terms and
$1/n!\leq1$. For $x=-t$, reciprocate
$1+t\leq\exp(t)\leq1/(1-t)$, using positivity.
\end{proof}
This argument supplies local bounds without assuming any continuity in
advance. The stronger inequality $\exp(x)\geq1+x$ for all real $x$ is
cited as homework in the source; these local bounds suffice here.

\begin{theorem}
The exponential function is continuous on $\mathbb R$.
\end{theorem}
\begin{proof}
As $x\to0$, both $1+x$ and $1/(1-x)$ tend to $1$. The sandwich theorem
gives $\exp(x)\to1=\exp(0)$.
For any fixed $c$,
\[
\exp(c+h)=\exp(c)\exp(h)\longrightarrow\exp(c)
\quad(h\to0),
\]
which proves continuity at $c$.
\end{proof}
For example, $\exp((x^3-7x)/(3x^2+5))$ is continuous on all of $\mathbb R$
by composition.

\section{Continuous functions on a closed bounded interval}
\begin{lemma}
If $x_n\in[a,b]$ and $x_n\to x$, then $x\in[a,b]$.
\end{lemma}
\begin{proof}
Pass to limits in $a\leq x_n\leq b$ using preservation of order under limits.
\end{proof}
Together with Bolzano--Weierstrass, this shows every sequence in $[a,b]$
has a subsequence converging to a point still in $[a,b]$. This property
is called \emph{sequential compactness}.

\begin{theorem}[Boundedness]
Every continuous function $f:[a,b]\to\mathbb R$ is bounded.
\end{theorem}
\begin{proof}
If it were unbounded, choose $x_n\in[a,b]$ with $|f(x_n)|>n$.
A subsequence $x_{j_n}$ converges to $c\in[a,b]$ by compactness.
Continuity implies $f(x_{j_n})\to f(c)$, hence this image subsequence is
bounded. But $|f(x_{j_n})|>j_n\geq n$, a contradiction.
\end{proof}

\begin{theorem}[Extreme value theorem]\label{thm:analysis-evt}
Every continuous function $f:[a,b]\to\mathbb R$ attains both a maximum
and a minimum: there exist $x_{\max},x_{\min}\in[a,b]$ with
\[
f(x_{\max})=\sup f([a,b]),\qquad
f(x_{\min})=\inf f([a,b]).
\]
\end{theorem}
\begin{proof}
Boundedness gives $M=\sup f([a,b])\in\mathbb R$.
For each $n$ choose $x_n\in[a,b]$ with $M-1/n<f(x_n)\leq M$.
A subsequence tends to $c\in[a,b]$; its image tends to $f(c)$ by continuity
and to $M$ by the sandwich theorem. Therefore $f(c)=M$.
Apply the same argument to $-f$ to obtain the minimum.
This completes the proof left unfinished on source p.~59.
\end{proof}

\begin{example}[Why the hypotheses matter]
The continuous function $1/x$ on $(0,1]$ is unbounded: the interval is not
closed. The continuous function $x$ on $[0,\infty)$ is unbounded: this
interval is closed but not bounded. Finally, on the compact interval $[0,1]$,
setting $f(0)=1$ and $f(x)=1/x$ for $x>0$ gives an unbounded function
because continuity at $0$ fails.
\end{example}

\section{The intermediate value theorem}
\begin{lemma}[Preservation of a strict inequality]
If $f$ is continuous at $c$ and $f(c)>\lambda$, then $f(x)>\lambda$
for all domain points sufficiently near $c$. An analogous statement holds
for $f(c)<\lambda$.
\end{lemma}
\begin{proof}
Choose $\varepsilon=(f(c)-\lambda)/2>0$ and require
$|f(x)-f(c)|<\varepsilon$. Then
$f(x)>f(c)-\varepsilon>\lambda$. Apply this to $-f$ for the other case.
\end{proof}

\begin{theorem}[Intermediate value theorem]\label{thm:analysis-ivt}
Let $f:[a,b]\to\mathbb R$ be continuous. Every real number $\lambda$
between $f(a)$ and $f(b)$ equals $f(c)$ for some $c\in[a,b]$.
\end{theorem}
\begin{proof}
The endpoint cases are immediate. Replacing $f$ by $-f$ if necessary,
assume $f(a)<\lambda<f(b)$. Set
$S=\{x\in[a,b]:f(x)<\lambda\}$ and $c=\sup S$.
Continuity at the endpoints gives points of $S$ strictly larger than $a$
and excludes a whole left neighbourhood of $b$ from $S$.
Consequently $a<c<b$.

If $f(c)<\lambda$, the lemma gives a point $x>c$ close to $c$ with
$f(x)<\lambda$, contradicting that $c$ bounds $S$.
If $f(c)>\lambda$, a neighbourhood $(c-\delta,c+\delta)$ has
$f(x)>\lambda$. Since $S$ contains no point above $c$ either,
$c-\delta$ would be an upper bound for $S$, after shrinking $\delta$
if necessary. This contradicts minimality of $c$.
Therefore $f(c)=\lambda$.
\end{proof}

Combining the extreme and intermediate value theorems gives
\[
f([a,b])=[m,M],\qquad m=\min f([a,b]),\quad M=\max f([a,b]).
\]
In particular a continuous function taking opposite signs at the endpoints
has a zero between them.

\section{The logarithm and real powers}
\begin{definition}[Bijection and inverse]
A function $f:A\to B$ is injective if different inputs give different
outputs, and surjective if every $y\in B$ has an input $x\in A$.
If both hold it is bijective and has an inverse $f^{-1}:B\to A$ with
$f^{-1}(f(x))=x$ and $f(f^{-1}(y))=y$.
\end{definition}

\begin{theorem}
The exponential is a strictly increasing bijection
$\exp:\mathbb R\to(0,\infty)$.
\end{theorem}
\begin{proof}
For $h>0$, its series gives $\exp(h)\geq1+h>1$. Thus if $y=x+h$,
\[
\exp(y)=\exp(x)\exp(h)>\exp(x).
\]
It is therefore injective. Also $\exp(x)\geq1+x\to+\infty$ as
$x\to+\infty$, and $\exp(x)=1/\exp(-x)\to0$ as $x\to-\infty$.
Given $u>0$, choose $x_1<x_2$ with $\exp(x_1)<u<\exp(x_2)$.
Continuity and the intermediate value theorem give $x$ with $\exp(x)=u$,
proving surjectivity.
\end{proof}

\begin{definition}[Natural logarithm]
The inverse of $\exp$ is the natural logarithm
$\log=\ln:(0,\infty)\to\mathbb R$. In particular,
\[
\exp(\log x)=x,\qquad \log(\exp t)=t.
\]
Define $e=\exp(1)=\sum_{n=0}^{\infty}1/n!\approx2.718281828459$.
\end{definition}

\begin{definition}[Real powers]
For $a>0$ and $x\in\mathbb R$, define $a^x=\exp(x\log a)$.
\end{definition}
This agrees with the usual definition for rational exponents: for positive
integers $q$, $(a^{p/q})^q=a^p$, and positivity selects the positive root.
The addition law gives, for $a>0$ and $x,y\in\mathbb R$,
\[
a^{x+y}=a^xa^y,\qquad (a^x)^y=a^{xy}.
\]
Also $e^x=\exp(x)$. If $a\ne1$, the map $x\mapsto a^x$ is a bijection
$\mathbb R\to(0,\infty)$, with inverse $\log_a$. For $a=1$ it is constant.
For a negative base, a real-valued $a^x$ cannot be defined for every real
$x$ with these same conventions.
